\documentclass[11pt]{article}
\usepackage[margin=1in]{geometry}

\usepackage{amsmath,amssymb}
\usepackage{amsthm}
\usepackage{booktabs}
\usepackage{multirow}
\usepackage{adjustbox}
\usepackage{graphicx}
\usepackage{algorithm}
\usepackage{algpseudocode}
\usepackage{natbib}
\usepackage{pifont}
\newcommand{\cmark}{\ding{51}}
\newcommand{\xmark}{\ding{55}}

\theoremstyle{plain}
\newtheorem{theorem}{Theorem}[section]
\newtheorem{lemma}[theorem]{Lemma}
\newtheorem{corollary}[theorem]{Corollary}
\theoremstyle{remark}
\newtheorem{remark}[theorem]{Remark}

\setlength{\parindent}{0pt}
\setlength{\parskip}{1\baselineskip}


\makeatletter
\newcommand{\affiliations}[1]{%
  \g@addto@macro\@author{\\[0.75em]
    \begin{minipage}{0.95\textwidth}
      \centering\normalsize #1
    \end{minipage}}%
}
\makeatother

\title{The Missing Signal: Gradient Injection Recovers $O(n^2)$ Residual Decoding in GNNs}
\author {
    % Authors
    Elad Shoham \textsuperscript{\rm 1},
    Havana Rika \textsuperscript{\rm 2},
    Dan Vilenchik\textsuperscript{\rm 1}
}
\affiliations {
    % Affiliations
    \textsuperscript{\rm 1}Ben-Gurion University of the Negev\\
    \textsuperscript{\rm 2} The Academic College of Tel Aviv-Yaffo \\
    shellad@post.bgu.ac.il, havanari@mta.ac.il, vilenchi@bgu.ac.il
}
\date{\today}

\theoremstyle{plain}
\newtheorem{claim}{Claim}[section]

\begin{document}
\maketitle

% ============================================================
\begin{abstract}
  Residual decoding algorithms score elements, prune one, and update state iteratively.
  Classical algorithms exploit additive statistics such as degree to perform each update
  in $O(n)$, achieving $O(n^2)$ total cost. Graph Nueral Networks (GNNs) cannot replicate this efficiency:
  nonlinear activations such as $\tanh$ or sigmoid make each unit of score change
  unequal---the gap between $\tanh(6)$ and $\tanh(5)$ differs fundamentally from that
  between $\tanh(2)$ and $\tanh(1)$---so no fixed additive correction exists after a
  deletion. This forces a full network rerun after every deletion, inflating cost to
  $O(n^3)$.

  We show that the objective gradient provides the missing signal and validate this on
  \textsc{Max-Clique} using planted-clique instances: after each vertex deletion, each
  surviving gradient coordinate changes by a local, maintainable quantity, updatable in
  $O(n)$ per step. Supplying this live gradient to a lightweight learned updater recovers
  $O(n^2)$ decoding, matching or exceeding the $O(n^3)$ rerun baseline across all
  difficulty regimes.
\end{abstract}

% ============================================================
\section{Introduction}
% ============================================================

A wide class of combinatorial algorithms proceeds by a \emph{residual loop}: score
the current elements, prune or select one, update the problem state, and repeat.
This loop underlies greedy peeling~\citep{feige2010}, spectral
rounding~\citep{alon1998}, branch-and-bound solvers, and local greedy methods more
broadly. Its tractability
rests on the \emph{additivity} of the maintained statistics. Degree is the
simplest example: when a vertex is removed, each survivor's degree changes by at
most one, depending only on local adjacency. This local, fixed correction makes the
full update $O(n)$, giving $O(n^2)$ total decoding on dense graphs.

\paragraph{Problem setting.}
We focus on \textsc{Max-Clique}: given a graph, find the largest subset of vertices
that are all pairwise adjacent. The decision version is NP-complete and the
optimization problem is NP-hard, making clique recovery a notorious testbed for
algorithmic heuristics. Classical decoders attack it with greedy top-$k$ selection,
local-degree refinement, spectral rounding, and related residual scoring rules. In
the planted-clique setting used in our experiments, the clique is hidden inside a
random background graph and the decoder must recover it. This makes
\textsc{Max-Clique} a natural place to ask whether learned scorers can replace these
hand-designed residual heuristics without losing their efficient update structure.

\paragraph{GNNs as learned scorers.}
GNNs have emerged as the natural candidate for replacing hand-crafted statistics
with learned scoring rules~\citep{karalias2020, min2022, garmendia2024}. Trained
end-to-end on problem instances, they produce vertex rankings for \textsc{Max-Clique},
\textsc{Max-Cut}, \textsc{Max-Independent-Set}, and related problems without
domain-specific feature engineering. Recent interpretability work shows that GNNs
trained on combinatorial tasks tend to rediscover classical algorithmic
heuristics~\citep{cappart2023}, and for \textsc{Max-Clique} specifically, degree-based
vertex ranking emerges consistently across architectures and training
regimes~\citep{shoham2026ltr}. This raises an optimistic picture: a GNN that has
internalized degree-based scoring should, in principle, slot directly into the
residual loop in place of the classical scorer.

\paragraph{The efficiency gap.}
This hope runs into a fundamental incompatibility. Classical residual algorithms work
because their maintained statistics update additively: removing a vertex changes each
survivor's degree by exactly one, so a single fixed correction suffices. GNN scores
cannot support this. Nonlinear activations such as $\tanh$ or sigmoid make each unit
of score change unequal---the gap between $\tanh(6)$ and $\tanh(5)$ differs from that
between $\tanh(2)$ and $\tanh(1)$---so no fixed additive correction exists after a
deletion. The problem compounds with scale: as instance size grows, score gaps between
adjacent candidates vanish, so even a small approximation error eventually inverts the
removal ordering. Once an inversion occurs, the approximate and exact decoders evolve
on structurally different residual graphs, and subsequent errors grow without bound.
The only safe remedy is to rerun the GNN after every deletion---$O(n^2)$ per step,
$O(n^3)$ total.

\paragraph{Central question.}
\textbf{\emph{Can a GNN be designed to support efficient incremental decoding,
    recovering the $O(n^2)$ complexity of classical residual algorithms without giving up
the expressive power of neural scoring?}}

\paragraph{The gradient as an additive coordinate.}
The obstruction above does not rule out efficient decoding; it rules out efficient
decoding from the \emph{raw GNN score} alone. The objective gradient provides exactly
the missing signal: after each vertex deletion, each surviving gradient coordinate
changes by a local, maintainable quantity determined entirely by the interaction
between the deleted vertex and each survivor. For the clique objective this correction
can be maintained in $O(n)$ per deletion by tracking three sufficient statistics,
giving $O(n^2)$ total cost after the initial GNN pass. Supplying this live cached
gradient to a learned score updater $U_\theta$ gives it access to the residual
correction directly, without requiring any inversion of the underlying nonlinear
activations.

\paragraph{Contributions.}
We identify a fundamental incompatibility between nonlinear activations and additive
residual updates: score gaps between adjacent candidates vanish with instance size,
so any fixed approximation error eventually inverts the removal ordering, after which
errors compound without bound. The objective gradient provides the missing signal,
as its per-deletion change is local and maintainable in $O(n)$ per step by tracking
three sufficient statistics, recovering $O(n^2)$ total decoding cost after the
initial GNN pass.

On the empirical side, we introduce U-LPR (Updatable Least Probable Removal), a
decoder that runs the GNN once and updates scores incrementally via a lightweight
learned updater $U_\theta$. We evaluate two variants: SGNN+U, which provides
$U_\theta$ with no gradient information, and SGNN+GU, which injects the cached
objective gradient into $U_\theta$ at each deletion step. We also introduce SIT
(Self-Imitation Training), a training regime that supervises $U_\theta$ by imitating
a rerun-based teacher. Both are described in full in Section~\ref{sec:method}.
Our findings:
\begin{itemize}
  \item \textbf{Recovery quality and running time.} SGNN+GU under SIT training
    achieves perfect approximation and exact recovery on Easy and Medium regimes
    at $O(n^2)$ cost, matching the $O(n^3)$ rerun baseline with a $4\times$
    wall-clock speedup.
  \item \textbf{Necessity of the gradient channel.} Removing the gradient from
    $U_\theta$ collapses U-LPR performance even under SIT training, confirming
    that gradient injection is a necessary architectural condition, not merely a
    helpful signal. Without it, SIT training makes performance strictly worse than
    objective training alone.
  \item \textbf{Learned residual concept.} Interpretability analysis reveals that
    the successful updater learns a coarse low-degree predicate: it avoids deleting
    vertices above the 15th percentile of residual degree among alive nodes, rather
    than sorting vertices by exact degree. The gradient channel is what makes this
    concept actionable during the residual loop.
\end{itemize}

\begin{figure}[t]
  \centering
  \includegraphics[width=0.92\textwidth]{figures/decoder_figure.png}
  \caption{Overview of the three neural decoder variants. SGNN uses the backbone
    scores directly, SGNN+U adds an incremental score updater without gradient
    information, and SGNN+GU supplies the updater with a cached residual gradient.
    The diagram highlights that the successful variant differs not only by adding an
    updater, but by exposing an update-compatible residual signal throughout the
    pruning trajectory.}
  \label{fig:decoder_overview}
\end{figure}

Figure~\ref{fig:decoder_overview} summarizes the architectural distinction between
SGNN, SGNN+U, and SGNN+GU: the gradient channel is the mechanism that makes the
learned updater compatible with residual decoding.

% ============================================================



\paragraph{Paper organization.}
Section~\ref{sec:related} reviews related work. Section~\ref{sec:method} develops
the methodology and introduces U-LPR and SIT. Section~\ref{sec:theory} presents
the analysis of why nonlinear scores fail and how the gradient resolves this.
Section~\ref{sec:experiments} reports the empirical evaluation.
Section~\ref{sec:concepts} analyzes the learned residual concept.
Section~\ref{sec:discussion} concludes with limitations and future directions.

% ============================================================
\section{Related Work}
\label{sec:related}
% ============================================================


\paragraph{Message-passing GNNs.}
Our backbone is a message-passing neural network (MPNN), which maintains a hidden
state $h_v^{(\ell)} \in \mathbb{R}^d$ updated over $L$ layers as
\[
  h_v^{(\ell+1)}
  =
  \sigma\!\left(
    \mathbf{M}\!\left[
      h_v^{(\ell)},\,
      \mathrm{Agg}_{u\in\mathcal{N}(v)} h_u^{(\ell)}
    \right]
  \right),
\]
with a per-vertex score $z_v = \mathbf{W}_{\mathrm{out}} h_v^{(L)}$ read off at the
end~\citep{xu2019}, optionally passed through a final nonlinearity such as sigmoid
to produce vertex probabilities $p_v \in (0,1)$. Their expressive power is bounded
by the 1-WL test~\citep{xu2019, morris2019}, though our results concern an
orthogonal limitation---\emph{update compatibility}---that persists even beyond the
WL bound. Since these networks ultimately output vertex scores through nonlinear
activations, a natural question is whether those scores resemble the heuristics used
by classical algorithms, and whether they can be updated as cheaply.

\paragraph{What GNNs learn on combinatorial problems.}
This score-based view connects directly to recent interpretability results.
\citet{cappart2023} argue that GNNs on NP-hard problems behave more like
polynomial-time heuristics than exact solvers, and \citet{shoham2026ltr} give a
concrete \textsc{Max-Clique} example: across architectures and datasets, learned
representations consistently recover degree-based vertex rankings. At first glance,
this is exactly what one would hope for, since degree is the canonical statistic
behind fast residual clique heuristics. The difficulty is that a useful residual
heuristic is not only a good ranking rule; it is also a rule whose state can be
updated cheaply after each deletion. Our work shows that this second property does
not automatically transfer from degree to a nonlinear GNN score.

\paragraph{GNNs for combinatorial optimization.}
This gap becomes central in neural combinatorial optimization, where learned scores
are routinely placed inside residual decoders. \citet{karalias2020} introduced an
unsupervised framework deriving GNN loss directly from the problem objective;
\citet{min2022} extended this with the scattering-based architecture we build on.
Similar unsupervised objectives have been applied to \textsc{Max-Cut},
\textsc{Max-Independent-Set}, and \textsc{TSP}~\citep{schuetz2022,garmendia2024}.
These works demonstrate strong combinatorial heuristics, but the standard decoding
pattern still reruns or reuses the network in a residual loop---precisely where
nonlinear activations make score changes unequal across the activation range, so
no fixed additive correction exists and accumulated errors eventually invert the
removal ordering.

\paragraph{Architectural alignment and gradient signals.}
Recent work on algorithm-aligned architectures shows that MPNNs can be made to
mirror classical algorithmic steps. \citet{yau2023} (OptGNN) proved that
polynomial-sized MPNNs can implement SDP gradient descent for Max-CSPs, obtaining
UGC-optimal guarantees in theory and outperforming SDP solvers in practice.
QRF-GNN~\citep{gaile2024} feeds intermediate GNN predictions back as dynamic node
features to improve convergence within a single inference pass, and \citet{hadou2024}
unroll a dual ascent algorithm into coupled networks so each layer computes a
meaningful algorithmic step.

Our work is complementary: rather than aligning message passing with a full
relaxation trajectory, we align the residual decoder with the local objective change
caused by each deletion. The algorithmic object we need is not a global optimization
dynamic but a live marginal signal that tracks how the objective shifts as the
residual graph shrinks, maintained in $O(n)$ per deletion via cached sufficient
statistics.

\paragraph{Imitation and self-supervised decoding.}
This training question connects our SIT regime to imitation learning in combinatorial
optimization~\citep{liu2021}, where a GNN is trained to replicate a stronger
algorithm's decisions. The novelty here is self-imitation: teacher and student share
weights, and the student must predict the rerun-based removal sequence from
incrementally updated scores using the cached gradient.

% ============================================================
\section{Methodology}
\label{sec:method}

\subsection{Problem Setting}
\label{sec:problem}

We focus on \textsc{Max-Clique}: given a graph $G=(V,E)$, find the largest subset
$C \subseteq V$ such that all pairs in $C$ are adjacent. We evaluate on the
planted-clique distribution $G(n,\tfrac{1}{2},k)$: an Erd\H{o}s--R\'enyi random
graph on $n$ vertices where each edge is included independently with probability
$1/2$, with a clique $C^\star$ of size $k$ planted by adding all missing edges among
$k$ uniformly chosen vertices. The goal is to recover a large valid clique, ideally
$C^\star$ itself.

\paragraph{Decoders.}
A natural family of decoders for this problem is the \emph{residual loop}: score all
vertices, remove the lowest-scoring one, and repeat until the remaining set forms a
clique, followed by an add-back pass that restores any removed vertex compatible with
the found clique. We evaluate five decoders:
\begin{itemize}
  \item \textbf{Top-$k$}: a non-iterative baseline that runs the GNN once, ranks
    vertices by score, and greedily adds them to the solution if adjacent to all
    previously selected vertices. $O(n^2)$.
  \item \textbf{LDR/LPR} (Least Degree/Probable Removal): the classical instantiation
    scores vertices by their current residual degree; LPR replaces degree with a
    learned probability $p_v$ produced by a GNN, retaining the same residual loop
    structure with scores fixed after the initial pass. Both are $O(n^2)$.
  \item \textbf{Rerun}: reruns the full GNN after every deletion to produce fresh
    scores. The only $O(n^3)$ decoder, serving as the adaptive upper bound.
  \item \textbf{U-LPR} (Updatable Least Probable Removal): runs the GNN once and
    applies a learned MLP updater $U_\theta$ after each deletion to incrementally
    correct the scores (Algorithm~\ref{alg:gub}). $O(n^2)$.
\end{itemize}


\paragraph{Datasets and their difficulty regimes.}
We evaluate on planted-clique instances $G(n,\tfrac{1}{2},k)$ with $n=1000$ vertices,
where $k$ is the size of the planted clique $C^\star$. Following~\citet{shoham2026ltr},
we partition instances by known theoretical thresholds on $k$:
\begin{itemize}
  \item \textbf{Easy} ($k\in[62,100]$): degree-based heuristics
    succeed~\citep{kucera1995}.
  \item \textbf{Medium} ($k\in[36,61]$): requires iterative residual
    methods~\citep{alon1998, feige2010}.
  \item \textbf{Hard} ($k\in[20,35]$): no polynomial-time algorithm is known to
    succeed~\citep{manurangsi2021}.
\end{itemize}

\subsection{Model Architecture}
\label{sec:architectures}

We evaluate three model families: SGNN, SGNN+U, and SGNN+GU. All share the same
message-passing backbone with $L=4$ steps, hidden dimension $64$, and $\sigma=\tanh$.
The input feature is a scalar liveness indicator; no degree, clique-membership, or
structural feature is provided. The backbone aggregates and updates as
\[
  q_v^{(\ell)} = \frac{1}{n-1}\sum_{u\in\mathcal{N}(v)} h_u^{(\ell)},
  \qquad
  h_v^{(\ell+1)} = \sigma\!\left(M_{\mathrm{upd}}\!\left[h_v^{(\ell)},
  q_v^{(\ell)}\right]\right),
\]
\[
  s_v = W_{\mathrm{out}} h_v^{(L)} + b_{\mathrm{out}},
  \qquad
  p_v = \operatorname{sigmoid}(s_v),
\]
where $M_{\mathrm{upd}}$ is a two-layer MLP and $p_v \in (0,1)$ is the vertex
probability used by the clique objective. SGNN uses this backbone alone. SGNN+U and
SGNN+GU add a lightweight MLP updater $U_\theta$ that corrects surviving scores after
each deletion of $v^*$:
\[
  s_i \leftarrow s_i + U_\theta(s_i,\, s_{v^*},\, A_{iv^*},\, \gamma_i,\,
  \gamma_{v^*}),
\]
where $\gamma_i, \gamma_{v^*}$ are gradient channels that differ between variants.
For SGNN+U, $\gamma_i = \gamma_{v^*} = 0$. For SGNN+GU,
\[
  \gamma_i = g_i = \frac{\partial \mathcal{L}_{\mathrm{clique}}}{\partial p_i},
  \qquad
  \gamma_{v^*} = g_{v^*} = \frac{\partial \mathcal{L}_{\mathrm{clique}}}{\partial
  p_{v^*}},
\]
maintained via cached sufficient statistics updated in $O(n)$ per deletion
(Section~\ref{sec:cached_maintenance}). \emph{The central question is whether this
gradient channel is necessary, sufficient, or both for efficient incremental decoding.}

\subsection{Cached Gradient Maintenance}
\label{sec:cached_maintenance}

Recomputing the full gradient after every deletion costs $O(n^2)$ per step and
$O(n^3)$ total. For our clique objective, this is unnecessary. The gradient at
each surviving vertex is a linear function of two cached sums: one aggregating
probabilities over neighbors, and one over non-neighbors. When a vertex $v$ is
deleted, its contribution to each sum is simply subtracted---an $O(n)$ operation
per deletion rather than a full $O(n^2)$ recomputation. The cache is initialized
in $O(n^2)$ alongside the initial GNN pass, so the total decoding cost remains
$O(n^2)$.


\subsection{Training}
\label{sec:training}

All models are trained on Medium-difficulty instances and evaluated across all three
regimes. Each run uses 64 epochs, with 1000 newly sampled instances per epoch,
minibatch size 32, and no overlap with the held-out evaluation instances. We note
that training on Easy or Hard instances yields identical results, confirming that
the findings are not sensitive to the choice of training regime. The unsupervised
clique objective is
\[
  \mathcal{L}_{\mathrm{clique}}(p; A, k)
  =
  -p^\top \mathbf{A} p
  +
  p^\top \mathbf{B} p
\]
where $p_i = \sigma(s_i)$ and $\mathbf{B}$ is the complement adjacency matrix.

\paragraph{Objective training (OBJ).}
$\mathcal{L}_{\mathrm{train}} = \mathcal{L}_{\mathrm{clique}}$. $U_\theta$ is
supervised only through the objective, with no direct signal on the removal sequence.

\paragraph{Self-Imitation Training (SIT).}
SIT adds an imitation loss that directly supervises the incremental decoder. At each
step $t$, a rerun-based teacher reruns the GNN on the current residual graph $G_t$
and selects the lowest-scoring vertex $v^t_{\mathrm{teach}}$. A student---running
the GNN only once and updating scores incrementally via $U_\theta$---is penalised
for disagreeing:
\[
  \mathcal{L}_{\mathrm{SIT}} = \frac{1}{T} \sum_{t=0}^{T-1}
  \mathrm{CE}_{\mathrm{softmax}}\!\left(-\tilde{s}^{(t)}_{V(G_t)},\;
  v^t_{\mathrm{teach}}\right),
  \qquad
  \mathcal{L}_{\mathrm{train}} = \mathcal{L}_{\mathrm{clique}} +
  \mathcal{L}_{\mathrm{SIT}}.
\]
For SGNN+U, $U_\theta$ receives no gradient signal and is trained purely through
backpropagation of $\mathcal{L}_{\mathrm{SIT}}$ with $\gamma_i = \gamma_{v^*} = 0$.
For SGNN+GU, the cached gradient is updated in $O(n)$ after each teacher removal
before being passed to $U_\theta$. In our experiments SIT uses $T=300$ teacher
removals per instance. The teacher cost is therefore paid at train time only, but
it remains the main training overhead.



\section{Claims and Proofs}
\label{sec:theory}

We formalize the three arguments sketched in the introduction.

\paragraph{Why nonlinear scores fail.}
Classical residual algorithms are efficient because their statistics update additively:
removing a vertex decreases each survivor's degree by exactly one, so a single fixed
correction suffices. For a GNN score to support the same property, it would need to
change by a fixed constant after every deletion.

\begin{claim}
  Exact additive updates require the score encoding to be affine in the maintained
  statistic. Nonlinear activations such as $\tanh$ or sigmoid are incompatible with
  this: the gap between $\tanh(6)$ and $\tanh(5)$ differs from the gap between
  $\tanh(2)$ and $\tanh(1)$, so no fixed additive correction exists after a deletion.
\end{claim}
\begin{proof}
  If $\varphi(d-1) = \varphi(d) - c$ for a fixed $c$ and all $d$, then every first
  discrete difference of $\varphi$ is constant, forcing $\varphi(d) = cd + b$. Any
  nonlinear activation $\sigma$ satisfies $\sigma(a-b) \neq \sigma(a) - \sigma(b)$
  in general, so no choice of parameters can make the post-deletion score equal the
  pre-deletion score minus a fixed constant.
\end{proof}

As instance size grows, score gaps between adjacent candidates shrink. This turns the
approximation error from a nuisance into a structural problem.

\begin{claim}
  Any fixed per-step approximation error $\varepsilon > 0$ eventually inverts the
  removal ordering. Once an inversion occurs, the approximate and exact decoders
  diverge onto structurally different residual graphs, after which errors compound
  without bound.
\end{claim}
\begin{proof}
  Let $\gamma_n$ be the minimum score gap. After $r = \lceil\gamma_n/\varepsilon\rceil$
  steps the accumulated error $r\varepsilon$ exceeds $\gamma_n$, inverting a pairwise
  ordering. After the inversion the two decoders remove different vertices; every
  subsequent score is evaluated on a different residual graph, so the discrepancy is
  no longer a perturbation of the original error but an uncontrolled structural
  difference.
\end{proof}

\paragraph{The gradient as the missing signal.}
The obstruction rules out efficient decoding from the raw score alone. The objective
gradient provides exactly what is missing.

\begin{claim}
  For any locally decomposable objective, deleting a vertex removes exactly its local
  contribution from each surviving gradient coordinate. For quadratic objectives this
  correction is linear in the deleted vertex's cached probability, and is directly
  representable in the first linear layer of $U_\theta$ before any nonlinearity fires.
\end{claim}
\begin{proof}
  Differentiating a locally decomposable objective, the gradient at vertex $i$
  involves only terms that include $i$. Deleting vertex $v$ removes exactly the
  pairwise term between $i$ and $v$, leaving all other terms unchanged. For a
  quadratic objective the removed term is linear in $p_v$, which is available as
  an input to $U_\theta$. The first affine layer of $U_\theta$ can therefore
  represent the correction exactly before any activation is applied. Without the
  gradient, recovering this correction from the nonlinear score $\sigma(z_v)$
  requires inverting $\sigma$, which the first claim shows is impossible with a
  fixed additive rule.
\end{proof}

\paragraph{Cached maintenance recovers $O(n^2)$.}

\begin{claim}
  The gradient cache can be initialized in $O(n^2)$ and updated in $O(n)$ per
  deletion, giving $O(n^2)$ total decoding cost after the initial GNN pass.
\end{claim}
\begin{proof}
  The gradient depends on two cached sums over neighbors and non-neighbors,
  each initialized by a matrix-vector product in $O(n^2)$. Deleting vertex $v$
  subtracts its contribution from each sum in $O(n)$. Over at most $n$ deletions
  the total update cost is $O(n^2)$, matching the initial pass.
\end{proof}


\section{Experiments}
\label{sec:experiments}

\subsection{Setup}
We use the architectures from Section~\ref{sec:architectures} and the training
protocol from Section~\ref{sec:training}. All experiments use planted-clique
instances with $n=1000$; evaluation set sizes are reported in the table captions.
Unless otherwise stated, reported values are mean $\pm$ standard deviation over
the corresponding held-out instances. Hardware: AMD Ryzen 9 5950X, 64\,GB RAM, NVIDIA RTX 3080;
Python 3.11, PyTorch 2.2.

\subsection{Main Results}
\label{sec:main_results}

Table~\ref{tab:main} reports planted clique recovery across all model and decoder
combinations.

\begin{table}[t]
  \centering
  \caption{Planted clique recovery on $n=1000$ across Easy, Medium, and Hard
    regimes, reported as mean$\,\pm\,$std over 10 held-out instances per regime;
    entries without $\pm$ are exactly $0.00$ or $1.00$ across all instances.
    \textbf{Bold} marks the highlighted best results: SGNN+GU+SIT with
    U-LPR is the strongest $O(n^2)$ incremental decoder on Easy and Medium, while
    SGNN+U+SIT with Rerun is best on Hard. Runtime reports asymptotic decoding cost:
    all decoders are $O(n^2)$ except Rerun, which is $O(n^3)$. SGNN has no updater
  and cannot run U-LPR.}
  \label{tab:main}
  \resizebox{\textwidth}{!}{%
    \begin{tabular}{lllccccccccc}
      \toprule
      & & & \multicolumn{3}{c}{Easy} & \multicolumn{3}{c}{Medium}
      & \multicolumn{3}{c}{Hard} \\
      \cmidrule(lr){4-6}\cmidrule(lr){7-9}\cmidrule(lr){10-12}
      Model & Decoder & Runtime & Approx & Overlap & Exact
      & Approx & Overlap & Exact
      & Approx & Overlap & Exact \\
      \midrule
      \multirow{3}{*}{SGNN + OBJ}
      & Top-$k$   & $O(n^2)$ & $1.00$ & $1.00$ & $1.00$
      & $0.75{\pm}0.43$ & $0.74{\pm}0.45$ & $0.67{\pm}0.58$
      & $0.46{\pm}0.12$ & $0.08{\pm}0.07$ & $0.00$ \\
      & Rerun     & $O(n^3)$ & $1.00$ & $1.00$ & $1.00$
      & $1.00$ & $1.00$ & $1.00$
      & $0.49{\pm}0.09$ & $0.00$ & $0.00$ \\
      & LPR       & $O(n^2)$ & $1.00$ & $1.00$ & $1.00$
      & $0.75{\pm}0.43$ & $0.74{\pm}0.45$ & $0.67{\pm}0.58$
      & $0.46{\pm}0.12$ & $0.08{\pm}0.07$ & $0.00$ \\
      \midrule
      \multirow{4}{*}{SGNN+U + OBJ}
      & Top-$k$   & $O(n^2)$ & $1.00$ & $1.00$ & $1.00$
      & $0.62{\pm}0.36$ & $0.58{\pm}0.40$ & $0.33{\pm}0.58$
      & $0.37{\pm}0.06$ & $0.08{\pm}0.07$ & $0.00$ \\
      & Rerun     & $O(n^3)$ & $1.00$ & $1.00$ & $1.00$
      & $1.00$ & $1.00$ & $1.00$
      & $0.63{\pm}0.33$ & $0.36{\pm}0.56$ & $0.33{\pm}0.58$ \\
      & LPR       & $O(n^2)$ & $1.00$ & $1.00$ & $1.00$
      & $0.62{\pm}0.36$ & $0.58{\pm}0.40$ & $0.33{\pm}0.58$
      & $0.37{\pm}0.06$ & $0.08{\pm}0.07$ & $0.00$ \\
      & U-LPR     & $O(n^2)$ & $0.13{\pm}0.01$ & $0.00$ & $0.00$
      & $0.17{\pm}0.01$ & $0.00$ & $0.00$
      & $0.34{\pm}0.09$ & $0.00$ & $0.00$ \\
      \midrule
      \multirow{4}{*}{SGNN+GU + OBJ}
      & Top-$k$   & $O(n^2)$ & $1.00$ & $1.00$ & $1.00$
      & $1.00$ & $1.00$ & $1.00$
      & $0.48{\pm}0.06$ & $0.06{\pm}0.03$ & $0.00$ \\
      & Rerun     & $O(n^3)$ & $1.00$ & $1.00$ & $1.00$
      & $1.00$ & $1.00$ & $1.00$
      & $0.54{\pm}0.08$ & $0.06{\pm}0.02$ & $0.00$ \\
      & LPR       & $O(n^2)$ & $1.00$ & $1.00$ & $1.00$
      & $1.00$ & $1.00$ & $1.00$
      & $0.50{\pm}0.05$ & $0.05{\pm}0.05$ & $0.00$ \\
      & U-LPR     & $O(n^2)$ & $0.12{\pm}0.00$ & $0.00{\pm}0.01$ & $0.00$
      & $0.18{\pm}0.02$ & $0.00$ & $0.00$
      & $0.42{\pm}0.04$ & $0.00$ & $0.00$ \\
      \midrule
      \multirow{4}{*}{SGNN+U + SIT}
      & Top-$k$   & $O(n^2)$ & $1.00$ & $1.00$ & $1.00$
      & $0.55{\pm}0.40$ & $0.48{\pm}0.47$ & $0.33{\pm}0.58$
      & $0.34{\pm}0.11$ & $0.04{\pm}0.01$ & $0.00$ \\
      & Rerun     & $O(n^3)$ & $1.00$ & $1.00$ & $1.00$
      & $1.00$ & $1.00$ & $1.00$
      & $\mathbf{0.82{\pm}0.31}$ & $\mathbf{0.67{\pm}0.58}$ & $\mathbf{0.67{\pm}0.58}$ \\
      & LPR       & $O(n^2)$ & $1.00$ & $1.00$ & $1.00$
      & $0.55{\pm}0.40$ & $0.48{\pm}0.47$ & $0.33{\pm}0.58$
      & $0.34{\pm}0.11$ & $0.04{\pm}0.01$ & $0.00$ \\
      & U-LPR     & $O(n^2)$ & $0.11{\pm}0.01$ & $0.00$ & $0.00$
      & $0.18{\pm}0.03$ & $0.00$ & $0.00$
      & $0.29{\pm}0.04$ & $0.00$ & $0.00$ \\
      \midrule
      \multirow{4}{*}{SGNN+GU + SIT}
      & Top-$k$   & $O(n^2)$ & $1.00$ & $1.00$ & $1.00$
      & $1.00$ & $1.00$ & $1.00$
      & $0.47{\pm}0.05$ & $0.02{\pm}0.03$ & $0.00$ \\
      & Rerun     & $O(n^3)$ & $1.00$ & $1.00$ & $1.00$
      & $1.00$ & $1.00$ & $1.00$
      & $0.50{\pm}0.07$ & $0.03{\pm}0.03$ & $0.00$ \\
      & LPR       & $O(n^2)$ & $1.00$ & $1.00$ & $1.00$
      & $0.46{\pm}0.21$ & $0.37{\pm}0.33$ & $0.00$
      & $0.47{\pm}0.05$ & $0.02{\pm}0.03$ & $0.00$ \\
      & U-LPR     & $O(n^2)$ & $\mathbf{1.00}$ & $\mathbf{1.00}$ & $\mathbf{1.00}$
      & $\mathbf{1.00}$ & $\mathbf{1.00}$ & $\mathbf{1.00}$
      & $0.48{\pm}0.09$ & $0.04{\pm}0.04$ & $0.00$ \\
      \bottomrule
    \end{tabular}%
  }
\end{table}

\paragraph{Rerun decoding is strong but expensive.}
Rerun decoding gives perfect Easy and Medium recovery for every model,
confirming that a strong GNN ranking suffices when fresh scores are available after
every deletion~\citep{shoham2026ltr}. Its cost is consistently higher:
Table~\ref{tab:timing} reports $0.76$--$0.78$\,s per instance for Rerun,
compared with $0.20$--$0.21$\,s for LPR and $0.37$--$0.57$\,s for U-LPR.

\paragraph{Without the gradient channel, U-LPR collapses.}
SGNN+U can still rank vertices reasonably under Top-$k$ and LPR, but its updater
does not work as a residual decoder. U-LPR stays near zero exact recovery: on Easy
it reaches only $0.11$--$0.13{\pm}0.01$ under different training regimes, while on
Medium it remains at $0.00$. Only Rerun restores perfect Easy/Medium recovery,
at the cost of recomputing the GNN after every deletion. SIT supervision alone is
therefore not enough without the gradient channel, and can even hurt U-LPR.

\paragraph{The gradient channel is necessary and SIT makes it usable.}
SGNN+GU under objective-only training still performs poorly under U-LPR
($0.12$--$0.18{\pm}0.02$ on Easy/Medium), showing that the gradient channel alone
is not enough---the model must also be trained to use it via SIT. Adding SIT
transforms U-LPR decoding to perfect approximation, overlap, and exact recovery on
Easy and Medium. The gradient channel makes the SIT objective achievable; SIT drives
the model to exploit the gradient signal at decoding time.

\paragraph{SIT trades static ranking for update quality.}
For SGNN+GU, SIT hurts LPR ($1.00\to0.46$ on Medium) while enabling perfect U-LPR
recovery ($1.00$). The model sacrifices one-pass score quality to learn an effective
incremental update. For SGNN+U this trade-off does not occur: without the gradient
channel, SIT cannot teach $U_\theta$ a useful update rule, so both LPR and U-LPR
remain poor.

\paragraph{Hard regime.}
Hard instances remain difficult for all models, consistent with the planted clique
hardness conjecture~\citep{manurangsi2021}. The strongest Hard result overall is
SGNN+U+SIT with Rerun decoding, which reaches $0.82{\pm}0.31$ approximation but
requires $O(n^3)$ time. Among $O(n^2)$ decoders, Hard is not solved reliably:
SGNN+GU+OBJ with LPR has the highest approximation at $0.50{\pm}0.05$, while
SGNN+GU+SIT with U-LPR is the strongest incremental variant at $0.48{\pm}0.09$.

\begin{table}[t]
  \centering
  \caption{Wall-clock time in seconds per instance on planted-clique graphs with
    $n=1000$, reported as mean $\pm$ standard deviation over 1000 held-out instances
  per regime and averaged over the Easy, Medium, and Hard regimes.}
  \label{tab:timing}
  \begin{tabular}{lcccc}
    \toprule
    Model & Top-$k$ & Rerun & LPR & U-LPR \\
    \midrule
    SGNN + OBJ      & 0.06 & 0.78 & 0.21 & --- \\
    SGNN+U + OBJ    & 0.06 & 0.77 & 0.20 & 0.38 \\
    SGNN+GU + OBJ   & 0.06 & 0.78 & 0.21 & 0.57 \\
    SGNN+U + SIT    & 0.06 & 0.76 & 0.20 & 0.37 \\
    SGNN+GU + SIT   & 0.06 & 0.76 & 0.20 & 0.54 \\
    \bottomrule
  \end{tabular}
\end{table}



\section{Interpretability and Concepts}
\label{sec:concepts}

Table~\ref{tab:pca} reports PCA diagnostics, Spearman correlations, and the
above-median deletion concept diagnostic. The goal is to separate what the GNN
backbone learns from what the residual updater uses during U-LPR.

\begin{table}[t]
  \centering
  \caption{Representation diagnostics on 1000 held-out test instances per regime
    ($n=1000$), reported as $|\rho|\pm\sigma$ (mean absolute Spearman $\pm$ std,
    all $p=0$ to machine precision unless noted).
    PC1/PC2: variance explained by the first two principal components of $h_v^{(L)}$.
    $|\rho|(\text{PC1},\text{deg})$ and $|\rho|(\text{logits},\text{deg})$: absolute
    Spearman between PC1 / output logits and residual degree.
    These four quantities are identical across all models, training regimes, and
    difficulty regimes: PC1$=1.00$, PC2$=0.00$,
    $|\rho|(\text{PC1},\text{deg})=1.00\pm0.00$,
    $|\rho|(\text{logits},\text{deg})=1.00\pm0.00$.
    $|\rho|(\tilde{s},\text{deg})$: absolute Spearman between post-update scores and
    residual degree after one application of $U_\theta$ (E$=$Easy, M$=$Medium,
    H$=$Hard). Despite similar global $|\rho|$ values at step~1, only SGNN+GU+SIT
    avoids removing above-15th-percentile-degree vertices:
    \emph{Above-p15} reports the fraction of U-LPR removal steps where the removed
    vertex has residual degree strictly above the 15th percentile among alive nodes
    (SIT training; medium regime, $n{=}1000$; $10$ graphs $\times$ 500 steps,
  pooled over all steps). SGNN has no updater and cannot run U-LPR.}
  \label{tab:pca}
  \resizebox{\textwidth}{!}{%
    \begin{tabular}{llcccccccccccccc}
      \toprule
      & & \multicolumn{2}{c}{PC var.} &
      \multicolumn{3}{c}{$|\rho|(\text{PC1},\text{deg})$} &
      \multicolumn{3}{c}{$|\rho|(\text{logits},\text{deg})$} &
      \multicolumn{3}{c}{$|\rho|(\tilde{s},\text{deg})$} &
      \multicolumn{2}{c}{Above-p15 (\%)} \\
      \cmidrule(lr){3-4}\cmidrule(lr){5-7}\cmidrule(lr){8-10}\cmidrule(lr){11-13}\cmidrule(lr){14-15}
      Model & Train & PC1 & PC2 & E & M & H & E & M & H & E & M & H & Pooled & Mean$\pm$std \\
      \midrule
      SGNN & OBJ
      & 1.00 & 0.00
      & $1.00\pm0.00$ & $1.00\pm0.00$ & $1.00\pm0.00$
      & $1.00\pm0.00$ & $1.00\pm0.00$ & $1.00\pm0.00$
      & \multicolumn{3}{c}{--- (no updater)}
      & \multicolumn{2}{c}{---} \\
      \midrule
      \multirow{2}{*}{SGNN+U}
      & OBJ
      & 1.00 & 0.00
      & $1.00\pm0.00$ & $1.00\pm0.00$ & $1.00\pm0.00$
      & $1.00\pm0.00$ & $1.00\pm0.00$ & $1.00\pm0.00$
      & $0.52\pm0.01$ & $0.52\pm0.01$ & $0.50\pm0.01$
      & --- & --- \\
      & SIT
      & 1.00 & 0.00
      & $1.00\pm0.00$ & $1.00\pm0.00$ & $1.00\pm0.00$
      & $1.00\pm0.00$ & $1.00\pm0.00$ & $1.00\pm0.00$
      & $0.50\pm0.04$ & $0.50\pm0.03$ & $0.49\pm0.02$
      & $85.3$ & $85.3\pm1.3$ \\
      \midrule
      \multirow{2}{*}{SGNN+GU}
      & OBJ
      & 1.00 & 0.00
      & $1.00\pm0.00$ & $1.00\pm0.00$ & $1.00\pm0.00$
      & $1.00\pm0.00$ & $1.00\pm0.00$ & $1.00\pm0.00$
      & $0.51\pm0.01$ & $0.47\pm0.03$ & $0.36\pm0.17$
      & --- & --- \\
      & SIT
      & 1.00 & 0.00
      & $1.00\pm0.00$ & $1.00\pm0.00$ & $1.00\pm0.00$
      & $1.00\pm0.00$ & $1.00\pm0.00$ & $1.00\pm0.00$
      & $0.49\pm0.03$ & $0.49\pm0.04$ & $0.51\pm0.02$
      & $\mathbf{0.0}$ & $\mathbf{0.0\pm0.0}$ \\
      \bottomrule
  \end{tabular}}
\end{table}

\paragraph{The GNN output is perfectly degree-ranked.}
The strongest signal is already present before any updater is applied: the final
GNN output logits have absolute Spearman correlation $1.00{\pm}0.00$ with vertex
degree across all models and regimes. Thus the backbone output learns a perfect
degree ranking, up to orientation, regardless of architecture or training regime.
The embeddings show the same collapse: PC1 explains essentially all variance
(PC1$=1.00$, PC2$=0.00$), matching the degree-collapse finding
of~\citet{shoham2026ltr}. Figure~\ref{fig:pca_degree_embeddings} shows this degree-aligned geometry for SGNN+U and SGNN+GU.

\begin{figure}[t]
  \centering
  \begin{minipage}{0.48\textwidth}
    \centering
    \includegraphics[width=\linewidth]{figures/SGNN_U_self_predictor_lpr_train_medium_g0_pc1_pc2_degree.pdf}

    \smallskip
    \textbf{(a)} SGNN+U.
  \end{minipage}\hfill
  \begin{minipage}{0.48\textwidth}
    \centering
    \includegraphics[width=\linewidth]{figures/SGNN_GU_self_predictor_lpr_train_medium_g0_pc1_pc2_degree.pdf}

    \smallskip
    \textbf{(b)} SGNN+GU.
  \end{minipage}
  \caption{Final-layer representation geometry. Each panel plots vertices in the
    first two principal-component coordinates of the last GNN layer, with points
    color-coded by residual degree. Both SGNN+U and SGNN+GU exhibit a strong
    alignment between degree and PC1, showing that the backbone representation is
    already organized by degree before the updater is applied.}
  \label{fig:pca_degree_embeddings}
\end{figure}


\paragraph{The updater learns a low-degree pruning concept.}
After one application of $U_\theta$, the absolute Spearman correlation between
residual degree and updated scores drops to $0.36$--$0.52$ across updater variants
and regimes, so the successful updater is not merely preserving the raw degree
ranking. The sharper diagnostic is the above-15th-percentile deletion rate:
SGNN+U+SIT removes an above-15th-percentile-degree vertex on $85.3\%$ of U-LPR
steps, while SGNN+GU+SIT removes one on $0.0\%$ of steps. Since SGNN+GU+SIT is
the only U-LPR variant that recovers Easy and Medium instances, the learned concept
is not exact degree sorting but a coarse low-degree predicate: avoid deleting
vertices above the 15th percentile of residual degree among alive nodes. The
gradient channel makes this concept actionable during the residual loop; without
it, the updater lacks the state needed to apply it reliably. Figures~\ref{fig:score_degree_pruning_t0},~\ref{fig:score_degree_pruning_t1},~\ref{fig:score_degree_pruning_t5},~\ref{fig:score_degree_pruning_t10}, and~\ref{fig:score_degree_pruning_t300} show the corresponding pruning decisions at iterations $t=0,1,5,10,300$.

\begin{figure}[t]
  \centering
  \begin{minipage}{0.48\textwidth}
    \centering
    \includegraphics[width=\linewidth]{figures/SGNN_U_self_predictor_lpr_train_medium_g0_upr_step0_score_vs_degree.pdf}

    \smallskip
    \textbf{(a)} SGNN+U, $t=0$.
  \end{minipage}\hfill
  \begin{minipage}{0.48\textwidth}
    \centering
    \includegraphics[width=\linewidth]{figures/SGNN_GU_self_predictor_lpr_train_medium_g0_upr_step0_score_vs_degree.pdf}

    \smallskip
    \textbf{(b)} SGNN+GU, $t=0$.
  \end{minipage}
  \caption{Score--degree pruning decision at U-LPR iteration $t=0$.
    Each panel plots updater score on the horizontal axis and residual degree on
    the vertical axis. The $\times$ marker denotes the vertex selected for removal,
    and the dashed horizontal line marks the 15th percentile of residual degree
    among alive vertices. Comparing SGNN+U with SGNN+GU shows whether the selected
    vertex lies below the low-degree threshold at this step.}
  \label{fig:score_degree_pruning_t0}
\end{figure}

\begin{figure}[t]
  \centering
  \begin{minipage}{0.48\textwidth}
    \centering
    \includegraphics[width=\linewidth]{figures/SGNN_U_self_predictor_lpr_train_medium_g0_upr_step1_score_vs_degree.pdf}

    \smallskip
    \textbf{(a)} SGNN+U, $t=1$.
  \end{minipage}\hfill
  \begin{minipage}{0.48\textwidth}
    \centering
    \includegraphics[width=\linewidth]{figures/SGNN_GU_self_predictor_lpr_train_medium_g0_upr_step1_score_vs_degree.pdf}

    \smallskip
    \textbf{(b)} SGNN+GU, $t=1$.
  \end{minipage}
  \caption{Score--degree pruning decision at U-LPR iteration $t=1$.
    Each panel plots updater score on the horizontal axis and residual degree on
    the vertical axis. The $\times$ marker denotes the vertex selected for removal,
    and the dashed horizontal line marks the 15th percentile of residual degree
    among alive vertices. Comparing SGNN+U with SGNN+GU shows whether the selected
    vertex lies below the low-degree threshold at this step.}
  \label{fig:score_degree_pruning_t1}
\end{figure}

\begin{figure}[t]
  \centering
  \begin{minipage}{0.48\textwidth}
    \centering
    \includegraphics[width=\linewidth]{figures/SGNN_U_self_predictor_lpr_train_medium_g0_upr_step5_score_vs_degree.pdf}

    \smallskip
    \textbf{(a)} SGNN+U, $t=5$.
  \end{minipage}\hfill
  \begin{minipage}{0.48\textwidth}
    \centering
    \includegraphics[width=\linewidth]{figures/SGNN_GU_self_predictor_lpr_train_medium_g0_upr_step5_score_vs_degree.pdf}

    \smallskip
    \textbf{(b)} SGNN+GU, $t=5$.
  \end{minipage}
  \caption{Score--degree pruning decision at U-LPR iteration $t=5$.
    Each panel plots updater score on the horizontal axis and residual degree on
    the vertical axis. The $\times$ marker denotes the vertex selected for removal,
    and the dashed horizontal line marks the 15th percentile of residual degree
    among alive vertices. Comparing SGNN+U with SGNN+GU shows whether the selected
    vertex lies below the low-degree threshold at this step.}
  \label{fig:score_degree_pruning_t5}
\end{figure}

\begin{figure}[t]
  \centering
  \begin{minipage}{0.48\textwidth}
    \centering
    \includegraphics[width=\linewidth]{figures/SGNN_U_self_predictor_lpr_train_medium_g0_upr_step10_score_vs_degree.pdf}

    \smallskip
    \textbf{(a)} SGNN+U, $t=10$.
  \end{minipage}\hfill
  \begin{minipage}{0.48\textwidth}
    \centering
    \includegraphics[width=\linewidth]{figures/SGNN_GU_self_predictor_lpr_train_medium_g0_upr_step10_score_vs_degree.pdf}

    \smallskip
    \textbf{(b)} SGNN+GU, $t=10$.
  \end{minipage}
  \caption{Score--degree pruning decision at U-LPR iteration $t=10$.
    Each panel plots updater score on the horizontal axis and residual degree on
    the vertical axis. The $\times$ marker denotes the vertex selected for removal,
    and the dashed horizontal line marks the 15th percentile of residual degree
    among alive vertices. Comparing SGNN+U with SGNN+GU shows whether the selected
    vertex lies below the low-degree threshold at this step.}
  \label{fig:score_degree_pruning_t10}
\end{figure}

\begin{figure}[t]
  \centering
  \begin{minipage}{0.48\textwidth}
    \centering
    \includegraphics[width=\linewidth]{figures/SGNN_U_self_predictor_lpr_train_medium_g0_upr_step300_score_vs_degree.pdf}

    \smallskip
    \textbf{(a)} SGNN+U, $t=300$.
  \end{minipage}\hfill
  \begin{minipage}{0.48\textwidth}
    \centering
    \includegraphics[width=\linewidth]{figures/SGNN_GU_self_predictor_lpr_train_medium_g0_upr_step300_score_vs_degree.pdf}

    \smallskip
    \textbf{(b)} SGNN+GU, $t=300$.
  \end{minipage}
  \caption{Score--degree pruning decision at U-LPR iteration $t=300$.
    Each panel plots updater score on the horizontal axis and residual degree on
    the vertical axis. The $\times$ marker denotes the vertex selected for removal,
    and the dashed horizontal line marks the 15th percentile of residual degree
    among alive vertices. Comparing SGNN+U with SGNN+GU shows whether the selected
    vertex lies below the low-degree threshold at this step.}
  \label{fig:score_degree_pruning_t300}
\end{figure}







\section{Discussion}
\label{sec:discussion}

The experiments point to a distinction that is easy to obscure when evaluating
learned combinatorial decoders: learning a good initial ranking is not the same
as learning a ranking that remains valid under residual updates. In our setting,
the GNN backbone reliably learns a degree-like ordering, and rerunning the GNN
after each deletion recovers strong Easy and Medium performance. Yet this does
not by itself yield an efficient decoder. The difficulty is not that the initial
score is uninformative; it is that the score is not intrinsically compatible with
the residual operation applied by the decoder.

This explains why the updater ablations are decisive. Without the gradient
channel, SIT gives the updater the right target sequence but not the right state
variable. The model is asked to imitate a rerun-based decoder from a nonlinear
score whose residual change is not directly available. This is why SGNN+U fails
under U-LPR even when trained with SIT. In contrast, SGNN+GU exposes a live
gradient-derived coordinate whose deletion effect is local and maintainable. SIT
then has a usable object to learn from: it does not need to invent the residual
correction from a static nonlinear score, but only to learn how to use the
maintained gradient signal to match the rerun decisions.

The results also clarify the role of cached maintenance. The cache is not merely
an implementation trick for speeding up the same computation. It changes what
information is available to the decoder throughout the pruning trajectory. A
stale gradient loses alignment because it ignores the vertices already removed;
a deletion-updated cache preserves the dominant residual change at $O(n)$ cost
per step. Thus the $O(n^2)$ decoder is not obtained by approximating an
$O(n^3)$ rerun more cheaply at each step, but by maintaining the right residual
quantity incrementally.

A useful consequence is that the successful updater need not preserve raw degree
alignment. The concept diagnostics show that the GNN score itself is degree-aligned,
while the post-update scores are only moderately correlated with degree. What
matters for U-LPR is the residual pruning concept: SGNN+GU+SIT avoids deleting
vertices above the 15th percentile of residual degree among alive nodes, whereas
SGNN+U does not. The gradient channel therefore supplies information beyond the
static degree ranking and lets the updater implement the low-degree deletion rule
needed for effective residual pruning.

More broadly, the paper suggests that neural decoders should be judged not only
by score quality but by update compatibility. Classical residual algorithms are
efficient because their maintained statistics have explicit local update rules.
Our results show that a learned decoder can recover this property by exposing
and maintaining a gradient coordinate. The problem-specific part is the form of
the objective gradient and its cached sufficient statistics; the general principle
is to align the learned decoder with the local objective change caused by each
deletion.

\paragraph{Limitations and future work.}
Our analysis suggests that gradient injection makes the residual correction
representable in the first linear layer of $U_\theta$, but it does not guarantee
that SGD will discover and use this consistently. GU addresses this by keeping
the gradient signal explicit during decoding. A stronger solution would be an
architecture or loss that makes the GNN itself learn update-compatible coordinates,
ideally removing the need for a separate updater. Our training procedure also relies
on SIT supervision from a rerun-based teacher, using $T=300$ teacher removals per
instance, which is expensive even though inference is quadratic. Reducing this
supervision budget, or replacing some teacher steps with cheaper targets, is an
important practical direction. Finally, the experiments focus on planted clique.
Extending the same gradient-maintenance principle to other locally decomposable
objectives, such as independent set and constraint-satisfaction relaxations, remains
an important test of the generality of the approach.



\bibliographystyle{plainnat}
\bibliography{bib}

% ============================================================
\clearpage
\appendix
\section{Algorithms}
\label{app:algorithms}
% ============================================================

\begin{algorithm}[t]
  \caption{U-LPR: Cached-Deletion Decoder}
  \label{alg:gub}
  \begin{algorithmic}[1]
    \State Run GNN once to obtain scores $s_i$ and probabilities $p_i$ for all $i$.
    \State Initialize gradient/cache from $p$. \hfill\Comment{$O(n^2)$}
    \State $C \leftarrow V(G)$;\; $H \leftarrow []$.
    \While{$C$ does not form a clique}
    \State $v^* \leftarrow \arg\min_{i \in C} s_i$;\; append $v^*$ to $H$;\; $C \leftarrow C \setminus \{v^*\}$.
    \State Remove $v^*$'s contribution from gradient cache. \hfill\Comment{$O(n)$}
    \State $g_i \leftarrow$ gradient from updated cache for all $i \in C$. \hfill\Comment{$O(n)$}
    \State $s_i \mathrel{+}= U_\theta(s_i,\, s_{v^*},\, A_{iv^*},\, g_i,\, g_{v^*})$ for all $i \in C$. \hfill\Comment{$O(n)$}
    \EndWhile
    \For{$v$ from last to first in $H$}
    \If{$C \cup \{v\}$ is a clique} $C \leftarrow C \cup \{v\}$. \EndIf
    \EndFor
    \State \Return $C$.
  \end{algorithmic}
\end{algorithm}

\begin{algorithm}[t]
  \caption{SIT Training Step}
  \label{alg:sit_training}
  \begin{algorithmic}[1]
    \State Sample instance $(G, k)$; run GNN once to obtain $\tilde{s}$, $p$.
    \State Initialize gradient/cache from $p$. \hfill\Comment{$O(n^2)$; SGNN+GU only}
    \For{$t = 0, \ldots, T-1$}
    \State $v^t_{\mathrm{teach}} \leftarrow \arg\min_{i \in V(G_t)} \mathrm{GNN}(G_t)_i$. \hfill\Comment{Teacher reruns GNN, $O(n^2)$}
    \State $\mathcal{L}_{\mathrm{SIT}} \mathrel{+}= \mathrm{CE}_{\mathrm{softmax}}(-\tilde{s}_{V(G_t)},\, v^t_{\mathrm{teach}})$. \hfill\Comment{Student predicts same removal}
    \State $G_t \leftarrow G_t \setminus \{v^t_{\mathrm{teach}}\}$.
    \State Update gradient cache by removing $v^t_{\mathrm{teach}}$. \hfill\Comment{$O(n)$; SGNN+GU only}
    \State $\tilde{s}_i \mathrel{+}= U_\theta(\tilde{s}_i,\, \tilde{s}_{v^t},\, A_{iv^t},\, \gamma_i,\, \gamma_{v^t})$ for all $i \in V(G_t)$.
    \EndFor
    \State Optimise $\mathcal{L}_{\mathrm{clique}} + \mathcal{L}_{\mathrm{SIT}} / T$.
  \end{algorithmic}
\end{algorithm}

\end{document}
